% =====================================================================
%  BAB 5: KONFIGURASI SISTEM & KUSTOMISASI LANJUTAN
%  (Settings, Users, Pods)
% =====================================================================

\chapter{Konfigurasi Sistem \& Kustomisasi Lanjutan}

Modul ini mencakup pengaturan global website, manajemen akun pengguna,
dan fitur \textit{Pods} (tipe konten kustom). Akses modul ini
\textbf{dibatasi hanya untuk Super Admin dan Admin}.

% -------------------------------------------------------------------
\section{Pengaturan Umum (\textit{General Settings})}
% -------------------------------------------------------------------

Pengaturan Umum menyimpan seluruh identitas sekolah yang ditampilkan di
berbagai halaman publik website.

\begin{figure}[H]
  \centering
  \includegraphics[width=0.9\textwidth]{../Img/Backend/screencapture-127-0-0-1-8000-admin-settings-2026-06-16-17_55_34.png}
  \caption{Tab Pengaturan Umum --- Identitas Sekolah}
  \label{fig:settings-general}
\end{figure}

\subsection{Panduan: Mengisi Identitas Sekolah}

\begin{enumerate}
  \item Klik menu \textbf{Pengaturan} (\textit{Settings}) di \textit{sidebar}.
        Anda langsung berada di tab \textbf{Umum}.
  \item Isi bagian \textbf{Identitas Sekolah}:
        \begin{itemize}
          \item \textbf{Nama Sekolah} --- nama resmi lengkap (wajib diisi)
          \item \textbf{Tagline} --- slogan atau motto sekolah
          \item \textbf{Deskripsi} --- paragraf perkenalan singkat
          \item \textbf{Alamat} --- alamat lengkap sekolah
          \item \textbf{Telepon} dan \textbf{Email}
          \item \textbf{Tahun Berdiri} dan \textbf{NPSN}
          \item \textbf{Akreditasi} dan \textbf{Nilai Akreditasi}
          \item \textbf{Status Sekolah} (Negeri/Swasta)
          \item \textbf{Luas Tanah (m2)}, \textbf{Luas Bangunan (m2)},
                \textbf{Waktu Belajar}, \textbf{Ruang Kelas}, \textbf{Program Studi}
        \end{itemize}
  \item Isi bagian \textbf{Logo \& Struktur}:
        \begin{itemize}
          \item \textbf{Logo Sekolah} --- digunakan juga di header, tab browser,
                dan favicon. Gunakan PNG berlatar transparan.
          \item \textbf{Gambar Struktur Organisasi} --- ditampilkan di halaman
                \texttt{/struktur-organisasi}.
        \end{itemize}
  \item Isi bagian \textbf{Profil Sekolah}:
        \begin{itemize}
          \item \textbf{Ringkasan Profil}, \textbf{Sejarah}, \textbf{Visi},
                \textbf{Misi} (baris baru per poin misi)
          \item \textbf{Video Profil (Link YouTube)} --- tempel URL penuh video,
                contoh: \texttt{https://www.youtube.com/watch?v=xxxxxxxx}.
                Video akan ditampilkan di beranda.
        \end{itemize}
  \item Isi bagian \textbf{Struktur \& Layanan}:
        \begin{itemize}
          \item \textbf{Deskripsi Struktur Organisasi}
          \item \textbf{Ringkasan Layanan} dan \textbf{Daftar Layanan}
                (baris baru per poin layanan)
        \end{itemize}
  \item Isi bagian \textbf{Statistik Beranda} --- angka ini ditampilkan di
        halaman depan sebagai \textit{stat counter}:
        \begin{itemize}
          \item \textbf{Total Guru \& Staf}
          \item \textbf{Siswa Aktif}
          \item \textbf{Program Unggulan}
          \item \textbf{Total Prestasi}
        \end{itemize}
  \item Isi bagian \textbf{Keamanan \& Akses}:
        \begin{itemize}
          \item \textbf{URL Login Admin Kustom} --- ganti \texttt{login} dengan
                kata rahasia Anda (misal: \texttt{portal-guru}).
                URL Login akan menjadi \texttt{/portal-guru}.
        \end{itemize}
  \item Klik tombol \textbf{Simpan Perubahan}.
\end{enumerate}

\begin{warnbox}[KRITIS: URL Login Kustom]
  Jika Anda mengubah \textbf{URL Login Admin Kustom} dan lupa kata rahasianya,
  Anda tidak akan bisa mengakses halaman login. Satu-satunya cara pemulihan adalah
  melalui \textit{database manager} (phpMyAdmin/Adminer) dengan mengubah nilai
  \texttt{admin\_login\_url} di tabel \texttt{settings} secara langsung.
  \textbf{Catat dan simpan URL rahasia ini di tempat yang aman!}
\end{warnbox}

% -------------------------------------------------------------------
\section{Pengaturan Kontak (\textit{Contact Settings})}
% -------------------------------------------------------------------

\begin{figure}[H]
  \centering
  \includegraphics[width=0.9\textwidth]{../Img/Backend/screencapture-127-0-0-1-8000-admin-settings-contact-2026-06-16-17_55_43.png}
  \caption{Tab Pengaturan Kontak --- WhatsApp dan Sosial Media}
  \label{fig:settings-contact}
\end{figure}

Navigasikan ke tab \textbf{Kontak} untuk mengisi:
\begin{itemize}
  \item \textbf{Maps Embed URL} --- URL \textit{iframe} Google Maps (untuk halaman \texttt{/kontak})
  \item \textbf{WhatsApp Number} --- nomor WhatsApp lengkap dengan kode negara (contoh: \texttt{628123456789})
  \item \textbf{Jam Operasional}
  \item Tautan media sosial: \textbf{Facebook URL}, \textbf{Instagram URL},
        \textbf{YouTube URL}, \textbf{Twitter/X URL}
\end{itemize}

Klik \textbf{Simpan Perubahan} setelah selesai.

% -------------------------------------------------------------------
\section{Pengaturan Tampilan (\textit{Appearance Settings})}
% -------------------------------------------------------------------

\begin{figure}[H]
  \centering
  \includegraphics[width=0.9\textwidth]{../Img/Backend/screencapture-127-0-0-1-8000-admin-settings-appearance-2026-06-16-17_55_58.png}
  \caption{Tab Pengaturan Tampilan --- Tema dan Tipografi}
  \label{fig:settings-appearance}
\end{figure}

Navigasikan ke tab \textbf{Tampilan} untuk mengatur:
\begin{itemize}
  \item \textbf{Warna Utama} dan \textbf{Warna Sekunder} --- kode warna HEX
        (contoh: \texttt{\#0ea5e9})
  \item \textbf{Jumlah Post per Halaman} --- berapa artikel yang ditampilkan
        per halaman di \texttt{/berita} (rentang: 1--60)
  \item \textbf{Font Utama (Teks Standar)} --- pilih dari daftar Google Fonts:
        \textit{Inter, Roboto, Poppins, Montserrat, Open Sans, Lato, Nunito,
        Outfit, Plus Jakarta Sans, Public Sans}
  \item \textbf{Font Judul (Aksen/Display)} --- pilih dari: \textit{Space Grotesk,
        Oswald, Bebas Neue, Raleway, DM Sans, Syne, Clash Display, Righteous, Cinzel}
  \item \textbf{Font Formal (Serif)} --- pilih dari: \textit{Playfair Display, Lora,
        Merriweather, PT Serif, Noto Serif, Crimson Text, EB Garamond}
\end{itemize}

Klik \textbf{Simpan Perubahan} setelah selesai.

% -------------------------------------------------------------------
\section{Pengaturan SEO}
% -------------------------------------------------------------------

\begin{figure}[H]
  \centering
  \includegraphics[width=0.9\textwidth]{../Img/Backend/screencapture-127-0-0-1-8000-admin-settings-seo-2026-06-16-17_56_11.png}
  \caption{Tab Pengaturan SEO Global}
  \label{fig:settings-seo}
\end{figure}

Navigasikan ke tab \textbf{SEO} untuk mengisi:
\begin{itemize}
  \item \textbf{Deskripsi Meta Default} --- deskripsi yang dipakai jika halaman
        tidak memiliki \textit{meta description} sendiri (maks. 320 karakter)
  \item \textbf{Google Analytics ID} --- kode pelacakan GA4
        (contoh: \texttt{G-XXXXXXXXXX})
  \item \textbf{Google Site Verification} --- kode verifikasi kepemilikan situs
        di Google Search Console
  \item \textbf{Robots.txt} --- isi aturan perayapan mesin pencari secara manual
\end{itemize}

% -------------------------------------------------------------------
\section{Pengaturan PPDB}
% -------------------------------------------------------------------

\begin{figure}[H]
  \centering
  \includegraphics[width=0.9\textwidth]{../Img/Backend/screencapture-127-0-0-1-8000-admin-settings-ppdb-2026-06-16-17_56_25.png}
  \caption{Tab Pengaturan PPDB --- Pendaftaran Peserta Didik Baru}
  \label{fig:settings-ppdb}
\end{figure}

Navigasikan ke tab \textbf{PPDB} untuk mengatur penerimaan siswa baru:
\begin{enumerate}
  \item Centang \textbf{Aktifkan Halaman PPDB} untuk memunculkan banner dan tombol
        ``Daftar PPDB'' di header website publik dan mengaktifkan halaman \texttt{/ppdb}.
  \item Isi \textbf{Link Pendaftaran (Opsional)} --- URL formulir pendaftaran eksternal
        (Google Form, website PPDB pusat, dll.).
  \item Isi \textbf{Tanggal Mulai (Opsional)} dan \textbf{Tanggal Berakhir (Opsional)}
        untuk menampilkan countdown atau informasi rentang waktu pendaftaran.
  \item Isi \textbf{Pengumuman / Informasi Tambahan PPDB} --- teks bebas untuk
        menampilkan syarat, narahubung, atau info gelombang pendaftaran.
  \item Klik tombol \textbf{Simpan Pengaturan}.
\end{enumerate}

% -------------------------------------------------------------------
\section{Manajemen Pengguna (\textit{Users})}
% -------------------------------------------------------------------

\begin{figure}[H]
  \centering
  \includegraphics[width=0.9\textwidth]{../Img/Backend/screencapture-127-0-0-1-8000-admin-users-2026-06-16-17_55_20.png}
  \caption{Daftar Pengguna Admin}
  \label{fig:users-index}
\end{figure}

\subsection{Panduan: Menambahkan Pengguna Baru}

\begin{enumerate}
  \item Klik menu \textbf{Pengguna} (\textit{Users}) di \textit{sidebar}.
  \item Klik tombol tambah.
  \item Isi \textbf{Nama} lengkap pengguna.
  \item Isi \textbf{Email} yang akan digunakan sebagai \textit{username} untuk login.
  \item Isi \textbf{Password} --- kata sandi sementara. Pengguna disarankan
        mengganti kata sandinya setelah login pertama kali.
        Saat mengedit pengguna yang sudah ada, \textbf{kosongkan kolom ini}
        jika tidak ingin mengubah kata sandi.
  \item Pilih \textbf{Role} dari \textit{dropdown}:
        \begin{itemize}
          \item \textit{Super Admin} --- akses penuh
          \item \textit{Admin} --- akses konten dan media
          \item \textit{Editor} --- akses tulis berita saja
        \end{itemize}
  \item Pastikan kotak \textbf{Aktif} dicentang agar pengguna dapat login.
  \item Klik tombol \textbf{Simpan}.
\end{enumerate}

\begin{notebox}[Menonaktifkan Pengguna]
  Untuk mencabut akses pengguna tanpa menghapus akunnya, \textbf{hapus centang}
  dari kotak \textbf{Aktif} dan simpan. Pengguna yang tidak aktif tidak dapat login
  ke panel admin.
\end{notebox}

% -------------------------------------------------------------------
\section{Tipe Konten Kustom (\textit{Pods})}
% -------------------------------------------------------------------

Fitur \textit{Pods} adalah mekanisme canggih yang memungkinkan Super Admin
membangun modul konten baru secara dinamis tanpa membutuhkan pemrograman.

\begin{figure}[H]
  \centering
  \includegraphics[width=0.9\textwidth]{../Img/Backend/screencapture-127-0-0-1-8000-admin-pods-2026-06-16-17_56_39.png}
  \caption{Halaman Manajemen Pods (Tipe Konten Kustom)}
  \label{fig:pods-index}
\end{figure}

\subsection{Skenario: Membuat Modul Testimoni Alumni}

Sekolah ingin menampilkan \textit{gallery} testimoni dari alumni terbaik di halaman depan.
Tidak ada modul bawaan untuk ini, maka Super Admin membuat \textit{Pod} baru.

\subsection{Panduan: Membuat Pod Baru}

\begin{enumerate}
  \item Buka menu \textbf{Pods} dan klik tambah.
  \item Isi \textbf{Nama Pod} (contoh: ``Testimoni Alumni'') dan deskripsinya.
  \item Simpan Pod. Sistem akan membuat struktur basis data secara otomatis.
  \item Klik \textbf{Kelola Field} pada Pod yang baru dibuat untuk mendefinisikan
        kolom-kolom datanya:
        \begin{itemize}
          \item Tambahkan field \textbf{Nama Lengkap} (tipe: Text)
          \item Tambahkan field \textbf{Angkatan} (tipe: Number)
          \item Tambahkan field \textbf{Kesan \& Pesan} (tipe: Textarea)
          \item Tambahkan field \textbf{Foto} (tipe: Image)
        \end{itemize}
  \item Setelah field tersimpan, buka menu \textbf{Pods $\rightarrow$ Testimoni Alumni $\rightarrow$ Entri}
        untuk mulai mengisi data testimoni alumni.
\end{enumerate}

\subsection{Menampilkan Data Pod di Konten Berita}

Data dari \textit{Pod} dapat disisipkan ke dalam konten berita menggunakan \textit{shortcode}
khusus di editor Quill:

\begin{lstlisting}
[pod:testimoni-alumni limit=5]
\end{lstlisting}

Sistem akan me-\textit{render} 5 entri testimoni terbaru dari Pod tersebut
secara dinamis di dalam konten artikel.

\subsection{Penanganan Masalah: Pods}

\begin{itemize}
  \item \textbf{\textit{Shortcode} tidak me-\textit{render}}:  Pastikan \textit{slug}
        Pod yang digunakan dalam \textit{shortcode} sama persis dengan \textit{slug}
        di sistem. Cek di halaman daftar Pods.
  \item \textbf{Data entri Pod tidak tersimpan}: Pastikan semua field yang
        didefinisikan sebagai \textit{required} sudah terisi sebelum menyimpan.
\end{itemize}
